
\documentclass[10pt]{article}
\usepackage[utf8]{inputenc}
\usepackage[T1]{fontenc}
\usepackage{amsmath, amssymb, xcolor, geometry, enumitem, multicol, tcolorbox}

\definecolor{mainblue}{RGB}{0, 102, 204}
\geometry{a4paper, margin=1.2cm, top=1cm, bottom=3cm, footskip=1.2cm}

\begin{document}
\enlargethispage{2.5cm}

% Modern Header
\begin{tcolorbox}[colback=mainblue!5, colframe=mainblue, sharp corners, boxrule=0.5pt]
    \small \textbf{Unit:} undefined \hfill \textbf{Name:} \rule{3cm}{0.4pt} \\
    \small \textbf{Topic:} Variables and Expressions / Order of Operations \hfill \textbf{Date:} \rule{2cm}{0.4pt} \\
    \centering \textbf{\large Homework 1}
\end{tcolorbox}

\vspace{0.2cm}

% MCQ Section
\begin{multicols}{2}
\begin{enumerate}[label=\textbf{Q\arabic*.}, leftmargin=*, itemsep=6pt, parsep=0pt]

  \item \small What is the variable in the expression $2x + 3$?
  \begin{enumerate}[label=(\Alph*), leftmargin=0.6cm, nosep, itemsep=2pt]
    \item $2$
    \item $3$
    \item $x$
    \item $2x$
    \item $2 + 3$
  \end{enumerate}

  \item \small Identify the coefficient in the term $4y$
  \begin{enumerate}[label=(\Alph*), leftmargin=0.6cm, nosep, itemsep=2pt]
    \item $y$
    \item $4$
    \item $4y$
    \item $1$
    \item $0$
  \end{enumerate}

  \item \small What is the constant term in the expression $2x - 5$?
  \begin{enumerate}[label=(\Alph*), leftmargin=0.6cm, nosep, itemsep=2pt]
    \item $2x$
    \item $-5$
    \item $2$
    \item $5$
    \item $-2x$
  \end{enumerate}

  \item \small In the expression $3 + 2$, what is the operator?
  \begin{enumerate}[label=(\Alph*), leftmargin=0.6cm, nosep, itemsep=2pt]
    \item $3$
    \item $2$
    \item $+$
    \item $-$
    \item $*$
  \end{enumerate}

  \item \small What does the term $2x^2$ represent?
  \begin{enumerate}[label=(\Alph*), leftmargin=0.6cm, nosep, itemsep=2pt]
    \item The coefficient of $x$ is $2$ and the exponent is $2$
    \item The variable is $2$ and the coefficient is $x^2$
    \item The coefficient of $x$ is $2$ and $x$ is squared
    \item The coefficient is $x$ and the variable is $2$
    \item The constant term is $2$ and the variable term is $x^2$
  \end{enumerate}

  \item \small In the expression $x + 2$, what operation is being performed?
  \begin{enumerate}[label=(\Alph*), leftmargin=0.6cm, nosep, itemsep=2pt]
    \item Subtraction
    \item Addition
    \item Multiplication
    \item Division
    \item Exponentiation
  \end{enumerate}

  \item \small Identify the term that is a variable in $3x + 2y - 4$
  \begin{enumerate}[label=(\Alph*), leftmargin=0.6cm, nosep, itemsep=2pt]
    \item $3x$ and $2y$
    \item $2y$ only
    \item $3x$ only
    \item $-4$
    \item $3x + 2y$
  \end{enumerate}

  \item \small What is the purpose of the order of operations?
  \begin{enumerate}[label=(\Alph*), leftmargin=0.6cm, nosep, itemsep=2pt]
    \item To simplify fractions
    \item To solve equations
    \item To evaluate expressions with multiple operations
    \item To graph functions
    \item To find the perimeter of a shape
  \end{enumerate}

\end{enumerate}
\end{multicols}

\vspace{-0.2cm}
\hrule
\vspace{0.2cm}
\textbf{\small Open Ended Questions (FRQ)}
\begin{enumerate}[label=\textbf{Q\arabic*.}, start=9, itemsep=5pt, topsep=0pt]

  \item \small Draw and explain how you would represent the expression $2x + 1$ using algebra tiles. Be sure to include a key for your drawing. \vspace{2cm}

  \item \small Explain, in your own words, what a variable is in an algebraic expression. Provide an example of a simple expression that contains a variable and describe what the variable represents in the context of the expression. \vspace{2cm}

\end{enumerate}

\vfill

% Chef's Tips Footer
\begin{tcolorbox}[colback=blue!5, colframe=mainblue!50, title=\small \textbf{Chef's Tips}, fonttitle=\bfseries, bottom=1mm]
\begin{itemize}[noitemsep, topsep=2pt, leftmargin=0.5cm]
  \item \scriptsize Focus on introducing basic concepts and definitions related to variables and expressions, such as coefficients and constants.\item \scriptsize Use simple, one-step problems to help students understand the fundamentals of the order of operations.\item \scriptsize Encourage students to use visual aids like algebra tiles to represent expressions and enhance their understanding of variables and expressions.
\end{itemize}
\end{tcolorbox}

\end{document}
  